% ============================================================
%  appendices/appendix-d-codebase.tex
%  Appendix D — Migration Codebase Reference
%  Full annotated listing of every file in the
%  nemotron-ascend-lifecycle project
% ============================================================

\chapter{Migration Codebase Reference}
\label{app:codebase}

This appendix provides the complete annotated source listing for
all files in the \texttt{nemotron-ascend-lifecycle} project
described in Chapter~4. Files are presented in execution order:
data preparation, operator implementations, training entry point,
configuration, inference, and tooling.

% ── D.1 ──────────────────────────────────────────────────────
\section{Project File Inventory}
\label{app:file-inventory}

\begin{table}[htbp]
  \centering
  \caption{Complete file inventory — \texttt{nemotron-ascend-lifecycle/}}
  \label{tab:file-inventory}
  \begin{tabularx}{\textwidth}{lXl}
    \toprule
    \tableheader{File & Purpose & Section} \\
    \midrule
    \texttt{env.toml}
      & Slurm cluster profiles and W\&B configuration
      & \secref{cluster-config} \\
    \rowcolor{TableRowOdd}
    \texttt{data/prepare\_datasets.py}
      & Download, filter, and format four HF datasets to JSONL
      & \secref{data-prep-pipeline} \\
    \texttt{ops/mamba\_ascend.py}
      & Mamba SSM selective scan — PyTorch fallback (sequential
        and chunked)
      & \secref{mamba-operator} \\
    \rowcolor{TableRowOdd}
    \texttt{ops/moe\_routing\_ascend.py}
      & MoE token dispatch, load balancing loss, dispatcher patch
      & \secref{moe-operator} \\
    \texttt{train\_sft\_ascend.py}
      & Patched SFT entry point — 6-step Ascend startup sequence
      & \secref{train-entry-point} \\
    \rowcolor{TableRowOdd}
    \texttt{configs/sft\_ascend.yaml}
      & Production SFT config — 8 NPUs, 3,400 iterations
      & \secref{sft-config} \\
    \texttt{configs/sft\_tiny.yaml}
      & Test config — 1 NPU, 100 iterations
      & \secref{test-config} \\
    \rowcolor{TableRowOdd}
    \texttt{inference/compile\_mindie.sh}
      & MindIE compilation from HF checkpoint
      & \secref{mindie-compilation} \\
    \texttt{inference/serve.py}
      & FastAPI inference server — MindIE/vLLM-Ascend backends
      & \secref{inference-server} \\
    \rowcolor{TableRowOdd}
    \texttt{inference/client.py}
      & Interactive client and test runner
      & \secref{inference-server} \\
    \texttt{inference/benchmark.py}
      & Async latency and throughput benchmark
      & \secref{perf-benchmark} \\
    \rowcolor{TableRowOdd}
    \texttt{tools/evaluate.py}
      & Technical QA evaluation — ROUGE-L, keyword coverage
      & \secref{qa-evaluation} \\
    \texttt{tools/validate\_loss.py}
      & W\&B loss curve comparison — CUDA vs.\ Ascend
      & \secref{loss-validation} \\
    \rowcolor{TableRowOdd}
    \path{docker/Dockerfile.sft}
      & SFT container — torch\_npu + Megatron-Bridge on Ascend
      & \secref{container-strategy} \\
    \path{docker/Dockerfile.inference}
      & Inference container — MindIE serving
      & \secref{container-strategy} \\
    \bottomrule
  \end{tabularx}
\end{table}

% ── D.2 ──────────────────────────────────────────────────────
\section{Dataset Formatters}
\label{app:dataset-formatters}

\tabref{formatter-mapping} documents how each raw dataset field
maps to the OpenAI chat format used for \acrshort{sft}.

\begin{table}[htbp]
  \centering
  \caption{Dataset field-to-chat-format mapping}
  \label{tab:formatter-mapping}
  \begin{tabularx}{\textwidth}{lXX}
    \toprule
    \tableheader{Dataset & Raw Fields & Chat Mapping} \\
    \midrule
    DialogStudio TechSupport
      & \texttt{dialog[].roles}, \texttt{dialog[].content}
      & \texttt{USR/USER} --\textgreater \texttt{user};
        \texttt{SYS/BOT} --\textgreater \texttt{assistant} \\
    \rowcolor{TableRowOdd}
    Arabic Alpaca
      & \texttt{instruction}, \texttt{input}, \texttt{output}
      Arabic Alpaca
        & \texttt{instruction + input} --\textgreater \texttt{user};
          \texttt{output} --\textgreater \texttt{assistant} \\
      Open-Platypus
        & \texttt{instruction + input} --\textgreater \texttt{user};
          \texttt{output} --\textgreater \texttt{assistant} \\
      \rowcolor{TableRowOdd}
      OPUS-100 fr-ar
        & \texttt{fr} --\textgreater \texttt{user};
          \texttt{ar} --\textgreater \texttt{assistant} \\
      \bottomrule
      \end{tabularx}
      \end{table}

      Quality filters applied per dataset:
      \begin{itemize}
      \item \textbf{DialogStudio}: minimum 2 turns; must have both
          \texttt{user} and \texttt{assistant} roles; first non-system
          turn must be \texttt{user}
      \item \textbf{Arabic Alpaca}: output length at least 20 characters
      \item \textbf{Open-Platypus}: at least one of 15 technical
          keywords in combined instruction + output text
      \item \textbf{OPUS-100}: French text at least 30 characters;
          at least one telecom keyword in French text
      \end{itemize}

% ── D.3 ──────────────────────────────────────────────────────
\section{Mamba SSM Technical Reference}
\label{app:mamba-reference}

The selective scan implements the input-dependent state recurrence 
across all timesteps. It computes the discretized state transition 
and input matrices before updating the hidden state and generating 
the output sequence.

\noindent Variable specifications for Nemotron 3 Nano:

\begin{table}[htbp]
  \centering
  \caption{Mamba SSM tensor dimensions}
  \label{tab:mamba-dims}
  \begin{tabularx}{0.85\textwidth}{lXX}
    \toprule
    \tableheader{Variable & Shape & Description} \\
    \midrule
    Input u     & (Batch, Dim, Len)     & Input sequence \\
    \rowcolor{TableRowOdd}
    Delta       & (Batch, Dim, Len)     & Time step (learned, input-dependent) \\
    Matrix A    & (Dim, N)              & State matrix (log-space, learned) \\
    \rowcolor{TableRowOdd}
    Matrix B    & (Batch, N, Len)       & Input projection (input-dependent) \\
    Matrix C    & (Batch, N, Len)       & Output projection (input-dependent) \\
    \rowcolor{TableRowOdd}
    Weight D    & (Dim,)                & Skip connection weight \\
    State x     & (Batch, Dim, N)       & Hidden state \\
    \rowcolor{TableRowOdd}
    Output y    & (Batch, Dim, Len)     & Output sequence \\
    \midrule
    Batch       & Batch size            & Typically 2--8 for SFT \\
    \rowcolor{TableRowOdd}
    Dim         & Model dimension       & 3,008 (Nemotron 3 Nano) \\
    N           & State size            & 16 (Mamba default) \\
    \rowcolor{TableRowOdd}
    Len         & Sequence length       & 4,096 (SFT pack size) \\
    \bottomrule
  \end{tabularx}
\end{table}

% ── D.4 ──────────────────────────────────────────────────────
\section{MoE Routing Technical Reference}
\label{app:moe-reference}

The routing computation selects the top-6 experts for each token 
based on router logits. The outputs of the selected experts are 
weighted by normalized routing probabilities and summed to produce 
the final MoE layer output.

\noindent For Nemotron 3 Nano: 128 total experts, 6 activated 
per token, model dimension 3,008, plus 2 shared experts 
which are always active.

The auxiliary load-balancing loss is added to the main 
cross-entropy loss during training to ensure uniform expert utilization.

% ── D.5 ──────────────────────────────────────────────────────
\section{CUDA-to-Ascend API Quick Reference}
\label{app:api-reference}

\begin{table}[htbp]
  \centering
  \caption{CUDA API to torch\_npu / AscendCL equivalent mapping}
  \label{tab:cuda-npu-api}
  \begin{tabularx}{\textwidth}{lXl}
    \toprule
    \tableheader{CUDA API & Ascend Equivalent & Change Type} \\
    \midrule
    \texttt{.cuda()}
      & \texttt{.npu()} / \texttt{transfer\_to\_npu()} auto-patch
      & Auto (monkey-patch) \\
    \rowcolor{TableRowOdd}
    \texttt{torch.cuda.is\_available()}
      & \texttt{torch.npu.is\_available()}
      & Manual if needed \\
    \texttt{torch.cuda.current\_device()}
      & \texttt{torch.npu.current\_device()}
      & Manual if needed \\
    \rowcolor{TableRowOdd}
    \texttt{torch.cuda.device\_count()}
      & \texttt{torch.npu.device\_count()}
      & Manual if needed \\
    \texttt{torch.cuda.amp.autocast}
      & \path{torch.npu.amp.autocast} via \path{_patch_autocast()}
      & Auto (patch in \path{train_sft_ascend.py}) \\
    \rowcolor{TableRowOdd}
    \texttt{torch.cuda.synchronize()}
      & \texttt{torch.npu.synchronize()}
      & Manual if needed \\
    \texttt{backend="nccl"}
      & \texttt{backend="hccl"} via env var
      & Config change \\
    \rowcolor{TableRowOdd}
    \texttt{NCCL\_DEBUG}
      & \path{HCCL_BUFFSIZE}, \path{HCCL_IF_BASE_PORT}
      & Env var rename \\
    \texttt{nvcc} kernel
      & TBE operator (\texttt{.py} DSL)
      & Full rewrite (future work) \\
    \rowcolor{TableRowOdd}
    \path{cudaMalloc}/\path{cudaFree}
      & \path{aclrtMalloc}/\path{aclrtFree} (AscendCL)
      & Not needed for PyTorch path \\
    \bottomrule
  \end{tabularx}
\end{table}

% ── D.6 ──────────────────────────────────────────────────────
\section{Environment Variable Reference}
\label{app:env-vars}

\begin{table}[htbp]
  \centering
  \caption{Ascend environment variables used in the migration}
  \label{tab:env-vars}
  \begin{tabularx}{\textwidth}{lXX}
    \toprule
    \tableheader{Variable & Description & Recommended Value} \\
    \midrule
    \path{ASCEND_VISIBLE_DEVICES}
      & Comma-separated NPU device IDs visible to the process
      & \texttt{"0,1,2,3,4,5,6,7"} \\
    \rowcolor{TableRowOdd}
    \path{HCCL_BUFFSIZE}
      & HCCL communication buffer size (MB)
      & \texttt{512} \\
    \path{HCCL_IF_BASE_PORT}
      & Base TCP port for HCCL inter-node sockets
      & \texttt{22118} \\
    \rowcolor{TableRowOdd}
    \path{HCCL_ALGO}
      & All-reduce algorithm (\texttt{ring} or \texttt{tree})
      & \texttt{ring} (multi-node) \\
    \path{ASCEND_GLOBAL_LOG_LEVEL}
      & CANN log verbosity (0=DEBUG, 1=INFO, 2=WARN, 3=ERROR)
      & \texttt{3} (production) / \texttt{0} (debug) \\
    \rowcolor{TableRowOdd}
    \path{USE_MAMBA_ASCEND_FALLBACK}
      & Enable PyTorch Mamba fallback (\texttt{1}/\texttt{0})
      & \texttt{1} \\
    \path{USE_MOE_ASCEND_FALLBACK}
      & Enable PyTorch MoE dispatch fallback (\texttt{1}/\texttt{0})
      & \texttt{1} \\
    \bottomrule
    \texttt{USE\_CHUNKED\_MAMBA}
      & Use chunked Mamba scan (faster) vs.\ sequential
      & \texttt{1} \\
    \texttt{TORCH\_DISTRIBUTED\_DEFAULT\_BACKEND}
      & Default \texttt{init\_process\_group} backend
      & \texttt{hccl} \\
    \bottomrule
  \end{tabularx}
\end{table}

% ── D.7 ──────────────────────────────────────────────────────
\section{Inference Server Request/Response Schema}
\label{app:api-schema}

\begin{lstlisting}[style=Python,
  caption={Chat completion request schema
           (\texttt{inference/serve.py})},
  label={lst:request-schema}]
# POST /v1/chat/completions
{
  "model":       "nemotron-nano3-telecom",   # required
  "messages": [
    {"role": "system",    "content": "..."},  # optional
    {"role": "user",      "content": "..."},  # required
    {"role": "assistant", "content": "..."},  # optional (multi-turn)
  ],
  "max_tokens":  512,     # default: 512
  "temperature": 0.7,     # default: 0.7
  "top_p":       0.9,     # default: 0.9
  "stream":      false,   # streaming not yet implemented
  "stop":        null     # optional stop sequences
}

# Response
{
  "id":      "chatcmpl-a1b2c3d4",
  "object":  "chat.completion",
  "created": 1717000000,
  "model":   "nemotron-nano3-telecom",
  "choices": [{
    "index":         0,
    "message":       {"role": "assistant", "content": "..."},
    "finish_reason": "stop"
  }],
  "usage": {
    "prompt_tokens":     142,
    "completion_tokens": 87,
    "total_tokens":      229
  },
  "_latency_ms": 1240.5    # server-side generation time
}
\end{lstlisting}

% ── D.8 ──────────────────────────────────────────────────────
\section{Benchmark Test Prompts}
\label{app:benchmark-prompts}

The benchmark in \filepath{inference/benchmark.py} cycles through
13 multilingual prompts of varying complexity. \tabref{benchmark-prompts}
documents all prompts and their expected complexity category.

\begin{table}[htbp]
  \centering
  \caption{Benchmark test prompts by language and complexity}
  \label{tab:benchmark-prompts}
  \begin{tabularx}{\textwidth}{llX}
    \toprule
    \tableheader{Lang & Complexity & Prompt} \\
    \midrule
    EN & Low   & ``What is the maximum data rate of 5G NR?'' \\
    \rowcolor{TableRowOdd}
    EN & Low   & ``How do I reset a Huawei OLT to factory defaults?'' \\
    EN & Low   & ``Explain the difference between OSPF and BGP.'' \\
    \rowcolor{TableRowOdd}
    EN & Low   & ``What does MTU stand for and what is its default value?'' \\
    EN & Medium & ``Intermittent packet loss on 10G optical link between
                   two CE6870 switches — diagnostic steps?'' \\
    \rowcolor{TableRowOdd}
    EN & Medium & ``Configure VLAN 100 with IP 192.168.100.1/24 on S5735
                   and enable DHCP.'' \\
    AR & Low   & \foreignlanguage{arabic}{``ما هو الفرق بين بروتوكول IPv4 و IPv6؟''} \\
    \rowcolor{TableRowOdd}
    AR & Low   & \foreignlanguage{arabic}{``كيف يمكنني تكوين OSPF على موجّه هواوي؟''} \\
    AR & Low   & \foreignlanguage{arabic}{``ما هي متطلبات الطاقة لمحطة قاعدة 5G NR؟''} \\
    \rowcolor{TableRowOdd}
    FR & Low   & ``Comment configurer un tunnel GRE sur un routeur Huawei?'' \\
    FR & Low   & ``Quelle est la différence entre la fibre monomode et multimode?'' \\
    \rowcolor{TableRowOdd}
    FR & Low   & ``Expliquez le processus de handover en 5G.'' \\
    EN & High  & ``50 Huawei AR6120 routers in hub-and-spoke MPLS topology,
                   more than 90 percent CPU on hub routers during business hours, 
                   BGP 500 routes, OSPF in access layer — causes and troubleshooting?'' \\
    \bottomrule
  \end{tabularx}
\end{table}
